% 参考文献
% 原书页码：354；标题
\BookUnnumbered{参考文献}
% 原书页码：354；标题
\BookSubUnnumbered{一、外文文献史料}
\par
\noindent\hangindent=2em\hangafter=1
% 原书页码：354；段落
[1] Dalai Lama V, The Illusive Play: The Autobiography of the Fifth Dalai Lama, Chi-%
cago: Serindia Publications, 2014.%
\par
\noindent\hangindent=2em\hangafter=1
% 原书页码：354；段落
[2] Donald St. Sure, Ray R. Noll and Edward Malatesta, eds., 100 Roman Documents%
\space%
Concerning the Chinese Rites Controversy (1645–1941), San Francisco: Ricci Insti-%
tute for Chinese-Western Cultural History, 1992.%
\par
\noindent\hangindent=2em\hangafter=1
% 原书页码：354；段落
[3] Ferdinand Verbiest, "Journeys into Tartary," in Pierre Joseph d'Orléans, History of%
\space%
the Two Tartar Conquerors of China. London: Hakluyt Society, 1854.%
\par
\noindent\hangindent=2em\hangafter=1
% 原书页码：354；段落
[4] Ferdinand Verbiest, The Astronomia Europaea of Ferdinand Verbiest, S. J. (Dillingen,%
\space%
1687): Text, Translation, Notes and Commentaries, Nettetal: Steyler Verlag, 1993.%
\par
\noindent\hangindent=2em\hangafter=1
% 原书页码：354；段落
[5] Joachim Bouvet, The History of Cang-Hy, the Present Emperor of China, London,%
\space%
printed for F. Coggan, 1699.%
\par
\noindent\hangindent=2em\hangafter=1
% 原书页码：354；段落
[6] John Bell, Travels from St. Petersburg, in Russia, to diverse parts of Asia, Glasgow,%
\space%
printed for the author by R. and A. Foulis, 1763.%
\par
\noindent\hangindent=2em\hangafter=1
% 原书页码：355；段落
[7] Kilian Stumpf, The Acta Pekinensia, or, Historical Records of the Maillard de Tournon%
\space%
Legation: First Transcribed Edition and English Annotated Translation, Rome: Institu-%
tum Historicum Societatis Iesu; Macau: Macau Ricci Institute, 2015.%
\par
\noindent\hangindent=2em\hangafter=1
% 原书页码：355；段落
[8] Kilian Stumpf, The Acta Pekinensia, or, Historical Records of the Maillard de Tour-%
non Legation. Volume II: September 1706–December 1707, Leiden: Brill, 2019.%
\par
\noindent\hangindent=2em\hangafter=1
% 原书页码：355；段落
[9] Ludovico Buglio, "An Abridgement of the Life and Death of F. Gabriel%
\space%
Magaillans," in Gabriel Magaillans, A New History of China, Containing a%
\space%
Description of the Most Considerable Particulars of That Vast Empire, London,%
\space%
printed for Thomas Newborough, at the Golden Ball, in S. Paul's Church-%
Yard, 1688.%
\par
\noindent\hangindent=2em\hangafter=1
% 原书页码：355；段落
[10] Matteo Ripa, Memoirs of Father Ripa During Thirteen Years' Residence at the Court%
\space%
of Peking in the Service of the Emperor of China, London: J. Murray, 1855.%
% 原书页码：355；标题
\BookSubUnnumbered{二、中文文献史料}
\par
\noindent\hangindent=2em\hangafter=1
% 原书页码：355；段落
[1]北京图书馆金石组编，《北京图书馆藏中国历代石刻拓本汇编》，郑州：%
中州古籍出版社，1989。%
\par
\noindent\hangindent=2em\hangafter=1
% 原书页码：355；段落
[2]北京行政学院著，《青石存史：“利玛窦和外国传教士墓地”的四百年沧%
桑》，北京：北京出版社，2011。%
\par
\noindent\hangindent=2em\hangafter=1
% 原书页码：355；段落
[3]《大清太宗文皇帝实录》，自《大清历朝实录》，台北：华文书局，1964。%
\par
\noindent\hangindent=2em\hangafter=1
% 原书页码：355；段落
[4]《圣祖仁皇帝实录》，自《大清历朝实录》，台北：华文书局，1964。%
\par
\noindent\hangindent=2em\hangafter=1
% 原书页码：355；段落
[5]《世宗宪皇帝实录》，自《大清历朝实录》，台北：华文书局，1964。%
\par
\noindent\hangindent=2em\hangafter=1
% 原书页码：355；段落
[6]《世祖章皇帝实录》，自《大清历朝实录》，台北：华文书局，1964。%
\par
\noindent\hangindent=2em\hangafter=1
% 原书页码：355；段落
[7]杜赫德编，《耶稣会士中国书简集：中国回忆录》，郑州：大象出版社，%
2001。%
\par
\noindent\hangindent=2em\hangafter=1
% 原书页码：356；段落
[8]杜文凯编，《清代西人见闻录》，北京：中国人民大学出版社，1985。%
\par
\noindent\hangindent=2em\hangafter=1
% 原书页码：356；段落
[9]故宫博物院明清档案部编，《关于江宁织造曹家档案史料》，北京：中华%
书局，1975。%
\par
\noindent\hangindent=2em\hangafter=1
% 原书页码：356；段落
[10]杭世骏撰，《武林理安寺志》，自《中国佛寺史志汇刊》，台北：明文书局，%
1980，第1辑，第21册。%
\par
\noindent\hangindent=2em\hangafter=1
% 原书页码：356；段落
[11]《皇朝文献通考》，自《景印文渊阁四库全书》，台北：台湾商务印书馆，%
1986。%
\par
\noindent\hangindent=2em\hangafter=1
% 原书页码：356；段落
[12]《皇清开国方略》，自《景印文渊阁四库全书》，台北：台湾商务印书馆，%
1986。%
\par
\noindent\hangindent=2em\hangafter=1
% 原书页码：356；段落
[13]《圣祖仁皇帝圣训》，自《景印文渊阁四库全书》，台北：台湾商务印书馆，%
1986。%
\par
\noindent\hangindent=2em\hangafter=1
% 原书页码：356；段落
[14]《圣祖仁皇帝庭训格言》，自《景印文渊阁四库全书》，台北：台湾商务%
印书馆，1986。%
\par
\noindent\hangindent=2em\hangafter=1
% 原书页码：356；段落
[15]陆世仪，《复社纪略》，上海：国粹学报馆，1908。%
\par
\noindent\hangindent=2em\hangafter=1
% 原书页码：356；段落
[16]《明熹宗实录》，台北：“中研院”历史语言研究所，1966。%
\par
\noindent\hangindent=2em\hangafter=1
% 原书页码：356；段落
[17]《钦定大清会典事例》，上海：商务印书馆，1908。%
\par
\noindent\hangindent=2em\hangafter=1
% 原书页码：356；段落
[18]佟世思，《先高曾祖三世行略》，自盛昱编，《八旗文经》，沈阳：辽沈书社，%
1988。%
\par
\noindent\hangindent=2em\hangafter=1
% 原书页码：356；段落
[19]王士禛，《池北偶谈》，北京：中华书局，1982。%
\par
\noindent\hangindent=2em\hangafter=1
% 原书页码：356；段落
[20] 叶梦珠，《阅世编》，北京：中华书局，2007。%
\par
\noindent\hangindent=2em\hangafter=1
% 原书页码：356；段落
[21]《雍亲王致年羹尧书》，自《文献丛编》，北京：故宫博物院，1930。%
\par
\noindent\hangindent=2em\hangafter=1
% 原书页码：356；段落
[22] 雍正编著，申国美整理，《御选语录》，北京：全国图书馆文献缩微复制%
中心，1993。%
\par
\noindent\hangindent=2em\hangafter=1
% 原书页码：356；段落
[23]《允禩允禟案：穆经远口供》，自《文献丛编》，北京：故宫博物院，1930。%
\par
\noindent\hangindent=2em\hangafter=1
% 原书页码：357；段落
[24]张伟仁主编，《明清档案》，台北：联经出版事业公司，1986。%
\par
\noindent\hangindent=2em\hangafter=1
% 原书页码：357；段落
[25] 赵尔巽等，《清史稿》，北京：中华书局，1976。%
\par
\noindent\hangindent=2em\hangafter=1
% 原书页码：357；段落
[26]昭梿，《啸亭杂录》，北京：中华书局，1980。%
\par
\noindent\hangindent=2em\hangafter=1
% 原书页码：357；段落
[27]中国第一历史档案馆编，《康熙朝汉文朱批奏折汇编》，北京：档案出版%
社，1984。%
\par
\noindent\hangindent=2em\hangafter=1
% 原书页码：357；段落
[28] 中国第一历史档案馆编，《康熙朝满文朱批奏折全译》，北京：中国社会%
科学出版社，1996。%
\par
\noindent\hangindent=2em\hangafter=1
% 原书页码：357；段落
[29] 中国第一历史档案馆编，《清中前期西洋天主教在华活动档案史料》，北%
京：中华书局，2003。%
\par
\noindent\hangindent=2em\hangafter=1
% 原书页码：357；段落
[30]中国第一历史档案馆、中国海外汉学研究中心合编，安双成编译，《清%
初西洋传教士满文档案译本》，郑州：大象出版社，2015。%
\par
\noindent\hangindent=2em\hangafter=1
% 原书页码：357；段落
[31] 中国第一历史档案馆，中国社会科学院历史研究所译注，《满文老档》，%
北京：中华书局，1990。%
\par
\noindent\hangindent=2em\hangafter=1
% 原书页码：357；段落
[32]“中研院”历史语言研究所编，《明清史料·乙编》，北京：国家图书馆%
出版社，2008。%
\par
\noindent\hangindent=2em\hangafter=1
% 原书页码：357；段落
[33]周岩编校，《熙朝定案》，自《明末清初天主教史文献新编》，北京：国%
家图书馆出版社，2013。%
\par
\noindent\hangindent=2em\hangafter=1
% 原书页码：357；段落
[34]周岩编校，《西洋新法历书·奏疏》，自《明末清初天主教史文献新编》，%
北京：国家图书馆出版社，2013。%
\par
\noindent\hangindent=2em\hangafter=1
% 原书页码：357；段落
[35] 周岩编校，《赠言》，自《明末清初天主教史文献新编》，北京：国家图%
书馆出版社，2013。%
% 原书页码：358；标题
\BookSubUnnumbered{三、专著及论文}
\par
\noindent\hangindent=2em\hangafter=1
% 原书页码：358；段落
[1] Aigle, Denise. The Mongol Empire between Myth and Reality: Studies in Anthropo-%
logical History. Leiden: Brill, 2014.%
\par
\noindent\hangindent=2em\hangafter=1
% 原书页码：358；段落
[2] Allan, C. W.. Jesuits at the Court of Peking. Shanghai: Kelly and Walsh, 1935.%
\par
\noindent\hangindent=2em\hangafter=1
% 原书页码：358；段落
[3] Allsop, Peter and Lindorff, Joyce. "Teodorico Pedrini: The Music and Letters of%
\space%
an 18th-century Missionary in China," in Vincentian Heritage Journal 27, no.2%
(2008).%
\par
\noindent\hangindent=2em\hangafter=1
% 原书页码：358；段落
[4]安双成，《汤若望案始末》，自《历史档案》，1992，第3期。%
\par
\noindent\hangindent=2em\hangafter=1
% 原书页码：358；段落
[5] Antonucci, Davor. "Notes on the Qing-Zunghar War (1690–1697): The Jesuits%
\space%
Sources," in Macau Ricci Institute, eds., Acta Pekinensia: Western Historical Sourc-%
es for the Kangxi Reign; International Symposium Organized by the Macau Ricci%
\space%
Institute, Macau, 5th–7th October, 2010. Macau: Macau Ricci Institute, 2013.%
\par
\noindent\hangindent=2em\hangafter=1
% 原书页码：358；段落
[6] Asen, Daniel. "Manchu Anatomy: Anatomical Knowledge and the Jesuits in%
\space%
Seventeenth- and Eighteenth-Century China," in Social History of Medicine 22,%
\space%
no.1 (2009).%
\par
\noindent\hangindent=2em\hangafter=1
% 原书页码：358；段落
[7] Bartlett, Beatrice. Monarchs and Ministers: The Grand Council in Mid-Ch'ing%
\space%
China, 1723-1820. Berkeley: University of California Press, 1990.%
\par
\noindent\hangindent=2em\hangafter=1
% 原书页码：358；段落
[8] Bayuk, Dimitri. "Father Antoine Gaubil, S. J. (1689-1759) and his Election to%
\space%
the Saint Petersburg Academy of Sciences," 2018, in \url{https://www.researchgate.net/publication/324953517}.%
\par
\noindent\hangindent=2em\hangafter=1
% 原书页码：358；段落
[9] Blussé, Leonard and Zurndorfer, Harriet T., eds., Conflict and Accommodation in%
\space%
Early Modern China: Essays in Honour of Erik Zürcher. Leiden: Brill, 1993.%
\par
\noindent\hangindent=2em\hangafter=1
% 原书页码：358；段落
[10] Brockey, Liam Matthew. Journey to the East: The Jesuit Mission to China 1579–%
1724. Cambridge: Harvard University Press, 2007.%
\par
\noindent\hangindent=2em\hangafter=1
% 原书页码：358；段落
[11] Cams, Mario. Companions in Geography: East-West Collaboration in the Mapping%
\space%
of Qing China (c. 1685-1735). Leiden: Brill, 2017.%
\par
\noindent\hangindent=2em\hangafter=1
% 原书页码：359；段落
[12] 曹树基，《中国人口史（第五卷清时期）》，上海：复旦大学出版社，2001。%
\par
\noindent\hangindent=2em\hangafter=1
% 原书页码：359；段落
[13] Chan, Albert. "A European Document on the Fall of the Ming Dynasty (1644-%
1649)," in Monumenta Serica 35 (1981-1983).%
\par
\noindent\hangindent=2em\hangafter=1
% 原书页码：359；段落
[14] Chang, Chia-feng. "Disease and its Impact on Politics, Diplomacy and the%
\space%
Military: The Case of Smallpox and the Manchus (1613-1795)," in Journal of the%
\space%
History of Medicine and Allied Sciences 57, no. 2 (2002).%
\par
\noindent\hangindent=2em\hangafter=1
% 原书页码：359；段落
[15]常江、李理，《清宫大内侍卫》，北京：故宫出版社，2013。%
\par
\noindent\hangindent=2em\hangafter=1
% 原书页码：359；段落
[16] Chang, Michael. A Court on Horseback: Imperial Touring \& the Construction of%
\space%
Qing Rule, 1680-1785. Cambridge: Harvard University Press, 2007.%
\par
\noindent\hangindent=2em\hangafter=1
% 原书页码：359；段落
[17] Chang, Ping-Ying. "Chinese Hereditary Mathematician Families of the Astro-%
nomical Bureau, 1620-1850." Ph. D. diss., The Graduate Center, City Universi-%
ty of New York, 2015.%
\par
\noindent\hangindent=2em\hangafter=1
% 原书页码：359；段落
[18] Chang, Te-chang, "The Economic Role of the Imperial Household in the Ching%
\space%
Dynasty," in The Journal of Asian Studies 31, no. 2 (1972).%
\par
\noindent\hangindent=2em\hangafter=1
% 原书页码：359；段落
[19] Chapman, Allan. "Tycho Brahe in China: The Jesuit Mission to Peking and the%
\space%
Iconography of European Instrument-making Processes," in Annals of Science 41,%
\space%
no. 5 (1984).%
\par
\noindent\hangindent=2em\hangafter=1
% 原书页码：359；段落
[20] 陈国栋，《清代内务府包衣三旗人员的分类及其旗下组织——兼论一些%
有关包衣的问题》，自《食货月刊》，1982，第12卷，第9期。%
\par
\noindent\hangindent=2em\hangafter=1
% 原书页码：359；段落
[21] 陈国栋，《康熙小臣养心殿总监造赵昌生平小考》，自冯明珠主编，《盛%
清社会与扬州研究》，台北：远流出版公司，2011。%
\par
\noindent\hangindent=2em\hangafter=1
% 原书页码：359；段落
[22] 陈佳华、傅克东，《八旗汉军考略》，自《民族研究》，1981，第5期。%
\par
\noindent\hangindent=2em\hangafter=1
% 原书页码：359；段落
[23] 陈寅恪，《柳如是别传》，北京：生活·读书·新知三联书店，2001，%
1980年初版。%
\par
\noindent\hangindent=2em\hangafter=1
% 原书页码：359；段落
[24]陈垣，《释奴才》，自《陈垣史学论著选》，上海：上海人民出版社，%
1981。%
\par
\noindent\hangindent=2em\hangafter=1
% 原书页码：359；段落
[25]陈垣，《汤若望与木陈忞》，自《陈垣学术论文集》，北京：中华书局，%
1980。%
\par
\noindent\hangindent=2em\hangafter=1
% 原书页码：360；段落
[26]陈垣，《雍乾间奉天主教之宗室》，自《陈垣学术论文集》，北京：中华%
书局，1980。%
\par
\noindent\hangindent=2em\hangafter=1
% 原书页码：360；段落
[27] Chu, Pingyi. "Scientific Dispute in the Imperial Court: The 1664 Calendar%
\space%
Case," in Chinese Science 14 (1997).%
\par
\noindent\hangindent=2em\hangafter=1
% 原书页码：360；段落
[28] Chu, T’ung-Tsu. Law and Society in Traditional China. Paris: Mouton, 1965.%
\par
\noindent\hangindent=2em\hangafter=1
% 原书页码：360；段落
[29] Cohen, Paul. China and Christianity: The Missionary Movement and the Growth of%
\space%
Chinese Antiforeignism 1860-1870. Cambridge: Harvard University Press, 1963.%
\par
\noindent\hangindent=2em\hangafter=1
% 原书页码：360；段落
[30] Collani, Claudia von. "Astronomy versus Astrology: Johann Adam Schall von%
\space%
Bell and his Superstitious' Chinese Calendar," in Archivum Historicum Societatis%
\space%
Iesu 82, no. 164.%
\par
\noindent\hangindent=2em\hangafter=1
% 原书页码：360；段落
[31]. "Kilian Stumpf and his Acta Pekinensia: Life, Context and Purpose of a%
\space%
Manuscript," in Macau Ricci Institute, eds., Acta Pekinensia: Western Historical%
\space%
Sources for the Kangxi Reign; International Symposium Organized by the Macau%
\space%
Ricci Institute, Macau, 5th-7th October, 2010. Macau: The Macau Ricci Institute,%
\space%
2013.%
\par
\noindent\hangindent=2em\hangafter=1
% 原书页码：360；段落
[32]. "Two Astronomers: Martino Martini and Johann Adam Schall von Bell,'%
in Luisa M. Paternicò, Collani, Claudia von and Scartezzini, Riccardo, eds.,%
\space%
Martino Martini: Man of Dialogue. Trento: Università di Trento, 2016.%
\par
\noindent\hangindent=2em\hangafter=1
% 原书页码：360；段落
[33]. "The Jesuit Rites Controversy," in Ines G. Županov, ed., The Oxford Hand-%
book of the Jesuits. New York: Oxford University Press, 2018.%
\par
\noindent\hangindent=2em\hangafter=1
% 原书页码：360；段落
[34]. "Charles Maigrot's Role in the Chinese Rites Controversy," in Mungello, D.%
\space%
E., ed., The Chinese Rites Controversy: Its History and Meaning. San Francisco: The%
\space%
Ricci Institute for Chinese-Western Cultural History, 1994.%
\par
\noindent\hangindent=2em\hangafter=1
% 原书页码：360；段落
[35] . "From the Earthly Court to the Heavenly Court," in Wardega and de%
\space%
Saldanhapp eds., In the Light and Shadow of an Emperor: Tomás Pereira, S.J.%
(1645–1708), the Kangxi Emperor and the Jesuit Mission in China. Cambridge:%
\space%
Cambridge Scholar Publishing, 2012.%
\par
\noindent\hangindent=2em\hangafter=1
% 原书页码：360–361；段落（跨页续段）
[36]. "Kangxi's Mandate of Heaven and Papal Authority, "in Deiwiks and Führer%
\space%
and Geulen, eds., Europe Meets China — China Meets Europe: The Beginnings of%
\space%
European-Chinese Scientific Exchange in the 17th Century. Routledge, 2014.%
\par
\noindent\hangindent=2em\hangafter=1
% 原书页码：361；段落
[37] . "Thomas and Tournon—Mission and Money," in Willy Vande Walle and%
\space%
Noel Golvers, eds., The History of the Relations between the Low Countries and%
\space%
China in the Qing Era (1644–1911). Leuven: Leuven University Press, 2003.%
\par
\noindent\hangindent=2em\hangafter=1
% 原书页码：361；段落
[38] Cooper, Michael. Rodrigues the Interpreter: An Early Jesuit in Japan and China.%
\space%
New York: Weatherhill, 1974.%
\par
\noindent\hangindent=2em\hangafter=1
% 原书页码：361；段落
[39] Crossley, Pamela. "Manchu Education," in Benjamin A. Elman and Alexander%
\space%
Woodside, eds., Education and Society in Late Imperial China, 1600–1900.%
\space%
Berkeley: University of California Press, 1994.%
\par
\noindent\hangindent=2em\hangafter=1
% 原书页码：361；段落
[40] . A Translucent Mirror: History and Identity in Qing Imperial Ideology.%
\space%
Berkeley: University of California Press, 1999.%
\par
\noindent\hangindent=2em\hangafter=1
% 原书页码：361；段落
[41] . "Slavery in Early Modern China," in David Eltis and Stanley L. Engerman,%
\space%
eds., The Cambridge World History of Slavery. Vol.3. Cambridge: Cambridge%
\space%
University Press, 2011.%
\par
\noindent\hangindent=2em\hangafter=1
% 原书页码：361；段落
[42] Dai, Yingcong. The Sichuan Frontier and Tibet: Imperial Strategy in the Early%
\space%
Qing. Seattle: University of Washington Press, 2011.%
\par
\noindent\hangindent=2em\hangafter=1
% 原书页码：361；段落
[43] D'Arelli, Francesco. "The Chinese College in Eighteenth-Century Naples," in%
\space%
East and West 58, no. 1/4 (2008).%
\par
\noindent\hangindent=2em\hangafter=1
% 原书页码：361；段落
[44] Dawson, Raymond. The Chinese Chameleon: An Analysis of European Conceptions%
\space%
of Chinese Civilisation. London: Oxford University Press, 1967.%
\par
\noindent\hangindent=2em\hangafter=1
% 原书页码：361；段落
[45] Deiwiks, Shu-Jyuan. "Some Cultural and Psychological Aspects of the Trial%
\space%
of Johann Adam Schall Before the Supreme Court of Peking According to the%
\space%
Secret Manchu Documents," in S. J. Deiwiks, B. Führer and T. Geulen, eds.,%
\space%
Europe Meets China—China Meets Europe: The Beginnings of European-Chinese%
\space%
Scientific Exchange in the 17th Century. Routledge, 2014.%
\par
\noindent\hangindent=2em\hangafter=1
% 原书页码：361–362；段落（跨页续段）
[46] Deiwiks, Shu-Jyuan. "The Secret Manchu Documents on the Trial of Jesuit Mis-%
sionary Johann Adam Schall (1592-1666) before the Supreme Court of Peking,%
\space%
in Monumenta Serica 51 (2003).%
\par
\noindent\hangindent=2em\hangafter=1
% 原书页码：362；段落
[47]董少新、黄一农，《崇祯年间招募葡兵新考》，自《历史研究》，2009，%
第5期。%
\par
\noindent\hangindent=2em\hangafter=1
% 原书页码：362；段落
[48]杜家骥，《清朝满蒙联姻研究》，北京：人民出版社，2003。%
\par
\noindent\hangindent=2em\hangafter=1
% 原书页码：362；段落
[49]杜家骥，《八旗与清朝政治论稿》，北京：人民出版社，2008。%
\par
\noindent\hangindent=2em\hangafter=1
% 原书页码：362；段落
[50]杜家骥，《清代八旗官制与行政》，北京：中国社会科学出版社，2015。%
\par
\noindent\hangindent=2em\hangafter=1
% 原书页码：362；段落
[51] 杜家骥，《清代内务府旗人复杂的旗籍及其多种身份——兼谈曹雪芹家%
族的旗籍及其身份》，自《民族研究》，2011，第3期。%
\par
\noindent\hangindent=2em\hangafter=1
% 原书页码：362；段落
[52] Dunne, George. Generation of Giants: The Story of the Jesuits in China in the last%
\space%
Decades of the Ming Dynasty. Notre Dame: University of Notre Dame Press, 1962.%
\par
\noindent\hangindent=2em\hangafter=1
% 原书页码：362；段落
[53] Elliott, Mark. The Manchu Way: The Eight Banners and Ethnic Identity in Late%
\space%
Imperial China. Stanford: Stanford University Press, 2001.%
\par
\noindent\hangindent=2em\hangafter=1
% 原书页码：362；段落
[54]. "Vocabulary Notes from the Manchu Archives 2: On the 'booi'," in Saksa-%
ha: A Journal of Manchu Studies 3 (1998).%
\par
\noindent\hangindent=2em\hangafter=1
% 原书页码：362；段落
[55]. "Ethnicity in the Qing Eight Banners," in Pamela Kyle Crossley, Helen%
\space%
Siu, and Donald Sutton, eds., Empire at the Margins: Culture, Ethnicity, and%
\space%
Frontier in Early Modern China. Berkeley and Los Angeles: University of%
\space%
California Press, 2006.%
\par
\noindent\hangindent=2em\hangafter=1
% 原书页码：362；段落
[56]. "The Manchu Language Archives and the Origins of the Palace Memorial%
\space%
System," in Late Imperial China 22 (2001).%
\par
\noindent\hangindent=2em\hangafter=1
% 原书页码：362；段落
[57] Elman, Benjamin. "New Perspectives of the Jesuits and Science in China: 1600-%
1800," in Bucknell Review 47, no. 2 (2004).%
\par
\noindent\hangindent=2em\hangafter=1
% 原书页码：362；段落
[58]. "The Jesuit Role as Experts' in High Qing Cartography and Technology," in%
\space%
Taida lishi xu bao 台大历史学报 31 (2003).%
\par
\noindent\hangindent=2em\hangafter=1
% 原书页码：362；段落
[59]. A Cultural History of Civil Examinations in Late Imperial China. Berkeley:%
\space%
University of California Press, 2000.%
\par
\noindent\hangindent=2em\hangafter=1
% 原书页码：362；段落
[60] Elverskog, Johan, Our Great Qing: The Mongols, Buddhism, and the State in Late%
\space%
Imperial China. Honolulu: University of Hawaii Press, 2008.%
\par
\noindent\hangindent=2em\hangafter=1
% 原书页码：363；段落
[61] Fairbank, J. K., and Teng, S. Y. "On the Ch'ing Tributary System," in Harvard%
\space%
Journal of Asiatic Studies 6, no. 2 (1941).%
\par
\noindent\hangindent=2em\hangafter=1
% 原书页码：363；段落
[62] 方豪，《清代旗人之信奉天主教与遭禁》，自《故宫文献》，1962，第4卷，%
第4期。%
\par
\noindent\hangindent=2em\hangafter=1
% 原书页码：363；段落
[63] 方豪，《中国天主教史人物传》，北京：中华书局，1988。%
\par
\noindent\hangindent=2em\hangafter=1
% 原书页码：363；段落
[64] Farquhar, David M. "Emperor as Bodhisattva in the Governance of the Ch’ing%
\space%
Empire," in Harvard Journal of Asiatic Studies 38, no. 1 (1978).%
\par
\noindent\hangindent=2em\hangafter=1
% 原书页码：363；段落
[65] Fatica, Michele. "Matteo Ripa Journal III (1716–1720) with New Documents%
\space%
for the History of the Kangxi Era," in Macau Ricci Institute, eds., in Acta%
\space%
Pekinensia: Western Historical Sources for the Kangxi Reign; International%
\space%
Symposium Organized by the Macau Ricci Institute, Macau, 5th-7th October 2010.%
\par
\noindent\hangindent=2em\hangafter=1
% 原书页码：363；段落
[66]冯尔康，《雍正传》，北京：人民出版社，1985。%
\par
\noindent\hangindent=2em\hangafter=1
% 原书页码：363；段落
[67] 冯明珠，《坚持与容忍——档案中所见康熙皇帝对中梵关系生变的因应》，%
自《中梵外交关系史国际学术研讨会论文集》，台北：辅仁大学出版社，%
2002。%
\par
\noindent\hangindent=2em\hangafter=1
% 原书页码：363；段落
[68] Feng, Jia. "The Emperor's Coffer: The Qing Imperial Fiscal Separation Between%
\space%
Privy Purse and State Treasury (1644–1912)." Ph. D. diss., UCLA, 2017.%
\par
\noindent\hangindent=2em\hangafter=1
% 原书页码：363；段落
[69]傅克东，《八旗户籍制度初探》，自《民族研究》，1983，第6期。%
\par
\noindent\hangindent=2em\hangafter=1
% 原书页码：363；段落
[70]傅克东，《从内佐领和管领谈到清代辛者库人》，自《清史研究通讯》，%
1986，第3期。%
\par
\noindent\hangindent=2em\hangafter=1
% 原书页码：363；段落
[71] Fu, Lo-shu. Documentary Chronicle of Sino-Western Relations (1644–1820).%
\space%
Tucson: The University of Arizona Press, 1966.%
\par
\noindent\hangindent=2em\hangafter=1
% 原书页码：363；段落
[72] 傅宗懋，《清初议政体制之研究》，自《政治大学学报》，1965，第11期。%
\par
\noindent\hangindent=2em\hangafter=1
% 原书页码：363；段落
[73] Givon, Yuval. "Connecting Eurasia: Jesuit Experimentation with Overland Mo-%
bility Between China and Europe, 1656–1664." Journal of World History 33, no.%
\space%
4 (2022).%
\par
\noindent\hangindent=2em\hangafter=1
% 原书页码：363；段落
[74] Gray, John Henry. Walks in the City of Canton. Hong Kong: De Souza \& Co.,%
\space%
1875.%
\par
\noindent\hangindent=2em\hangafter=1
% 原书页码：364；段落
[75] Gernet, Jacques. China and the Christian Impact: A Conflict of Culture. Cam-%
bridge: Cambridge University Press, 1982.%
\par
\noindent\hangindent=2em\hangafter=1
% 原书页码：364；段落
[76] 古洛东，《圣教人川记》，成都：四川人民出版社，1981，1918初版。%
\par
\noindent\hangindent=2em\hangafter=1
% 原书页码：364；段落
[77] Guy, R. Kent. "Who Were the Manchus? A Review Essay," in Journal of Asian%
\space%
Studies 61, no.1 (2002).%
\par
\noindent\hangindent=2em\hangafter=1
% 原书页码：364；段落
[78] 韩琦，《姗姗来迟的“西洋消息”——1709年教皇致康熙信到达宫廷始%
末》，自《文化杂志》，2005，第55期。%
\par
\noindent\hangindent=2em\hangafter=1
% 原书页码：364；段落
[79]韩琦、吴旻校注，《熙朝崇正集熙朝定案（外三种）》，北京：中华书局，%
2006。%
\par
\noindent\hangindent=2em\hangafter=1
% 原书页码：364；段落
[80] Hanson, Marta and Gianna Pomata, "Travels of a Chinese Pulse Treatise: The%
\space%
Latin and French Translations of the Tuzhu maijue bianzhen 圖註脈訣辨真%
(1650s-1730s)," in Harold J. Cook, ed., Translation at Work: Chinese Medicine in%
\space%
the First Global Age, Leiden: Brill, 2020.%
\par
\noindent\hangindent=2em\hangafter=1
% 原书页码：364；段落
[81] Hevia, James. "Lamas, Emperors, and Rituals: Political Implications in Qing%
\space%
Imperial Ceremonies," in Journal of the International Association of Buddhist%
\space%
Studies 16 (1993).%
\par
\noindent\hangindent=2em\hangafter=1
% 原书页码：364；段落
[82] Hsia, Florence. "Some Observations on the Observations: The Decline of the%
\space%
French Jesuit Scientific Mission in China," in Revue de Synthèse 120 (1999).%
\par
\noindent\hangindent=2em\hangafter=1
% 原书页码：364；段落
[83] . Sojourners in a Strange Land: Jesuits and Their Scientific Missions in Late%
\space%
Imperial China. Chicago: University of Chicago Press, 2009.%
\par
\noindent\hangindent=2em\hangafter=1
% 原书页码：364；段落
[84] Hsia, R. Po-Chia. "Imperial China and Jesuit Mission," in Roberto Ribeiro and%
\space%
John O'Malley, eds., Jesuit Mapmaking in China. Philadelphia: Saint Joseph's%
\space%
University Press, 2014.%
\par
\noindent\hangindent=2em\hangafter=1
% 原书页码：364；段落
[85] . "Jesuit Survival and Restoration in China," in Robert A. Maryks and%
\space%
Jonathan Wright, eds., Jesuit Survival and Restoration: A Global History, 1773–%
1900. Leiden: Brill, 2014.%
\par
\noindent\hangindent=2em\hangafter=1
% 原书页码：364–365；段落（跨页续段）
[86]. "The End of the Jesuit Mission in China," in Jeffrey D. Burson and%
\space%
Jonathan Wright, eds., The Jesuit Suppression in Global Context: Causes, Events, and%
\space%
Consequences. Cambridge: Cambridge University Press, 2015.%
\par
\noindent\hangindent=2em\hangafter=1
% 原书页码：365；段落
[87] Hu, Xiangyu. "The Juridical System of the Qing Dynasty in Beijing, 1644–1900."%
Ph.D. diss., University of Minnesota, 2011.%
\par
\noindent\hangindent=2em\hangafter=1
% 原书页码：365；段落
[88]黄伯禄，《正教奉褒》，上海：慈母堂，1904。%
\par
\noindent\hangindent=2em\hangafter=1
% 原书页码：365；段落
[89] Huang, Pei. Autocracy at Work: A Study of the Yung-cheng Period: 1723–1735.%
\space%
Bloomington: Indiana University Press, 1976.%
\par
\noindent\hangindent=2em\hangafter=1
% 原书页码：365；段落
[90] Huang, Ray. 1587, The Year of No Significance: The Ming Dynasty in Decline. New%
\space%
Haven: Yale University Press, 1981.%
\par
\noindent\hangindent=2em\hangafter=1
% 原书页码：365；段落
[91] 黄一农，《清前期对觜、参两宿先后次序的争执：社会天文学史之一个案%
研究》，自杨翠华，黄一农主编，《近代中国科技史论集》，台北：“中研院”%
近代史研究所，1991。%
\par
\noindent\hangindent=2em\hangafter=1
% 原书页码：365；段落
[92]黄一农，《择日之争与“康熙历狱”》，自《清华学报》，1991，第21卷，%
第2期。%
\par
\noindent\hangindent=2em\hangafter=1
% 原书页码：365；段落
[93] 黄一农，《康熙朝汉人士大夫对“历狱”的态度及其所衍生的传说》，《汉%
学研究》，1993，第11卷，第2期。%
\par
\noindent\hangindent=2em\hangafter=1
% 原书页码：365；段落
[94]黄一农，《耶稣会士对中国传统星占术数的态度》，自《九州学刊》，%
1991，第4卷，第3期。%
\par
\noindent\hangindent=2em\hangafter=1
% 原书页码：365；段落
[95] 黄一农，《张宸生平及其与杨光先间的冲突》，自《九州学刊》，1993，第6%
卷，第1期。%
\par
\noindent\hangindent=2em\hangafter=1
% 原书页码：365；段落
[96]黄一农，《红夷大炮与皇太极创立的八旗汉军》，自《历史研究》，2004，%
第4期。%
\par
\noindent\hangindent=2em\hangafter=1
% 原书页码：365；段落
[97]黄一农，《天主教徒孙元化与明末传华的西洋火炮》，自《“中研院”历%
史语言研究所集刊》，1996，第67卷，第4期。%
\par
\noindent\hangindent=2em\hangafter=1
% 原书页码：365；段落
[98] Hucker, Charles. The Censorial System of Ming China. Stanford: Stanford Univer-%
sity Press, 1966.%
\par
\noindent\hangindent=2em\hangafter=1
% 原书页码：365；段落
[99] Hughes, Gerard. "The Acta's Four Portraits," in Macau Ricci Institute, eds., Acta%
\space%
Pekinensia: Western Historical Sources for the Kangxi Reign. Macau: Macau Ricci%
\space%
Institute, 2013.%
\par
\noindent\hangindent=2em\hangafter=1
% 原书页码：365；段落
[100] Hummel, Arthur, ed., Eminent Chinese of the Ch’ing Period (1644–1912). Wash-%
ington: U.S. Government Printing Office, 1944.%
\par
\noindent\hangindent=2em\hangafter=1
% 原书页码：366；段落
[101] Jaskov, Helena. "The Negotiated Geography of the Treaty of Nerchinsk (1689)%
\space%
and the Role of the Jesuits," in Late Imperial China 40, no. 2 (2019).%
\par
\noindent\hangindent=2em\hangafter=1
% 原书页码：366；段落
[102] Jami, Catherine. "Tomé Pereira (1645-1708); Clockmaker; Musician and In-%
terpreter at the Kangxi Court: Portuguese Interests and the Transmission of Sci-%
ence," in Luis Saraiva and Catherine Jami, eds., The Jesuits, the Padroado and East%
\space%
Asian Science (1552–1773). Singapore: World Scientific Press, 2008.%
\par
\noindent\hangindent=2em\hangafter=1
% 原书页码：366；段落
[103]. "Revisiting the Calendar Case (1664–1669): Science, Religion, and Poli-%
tics in Early Qing Beijing," in Korean Journal of History of Science 27, no. 2 (2015).%
\par
\noindent\hangindent=2em\hangafter=1
% 原书页码：366；段落
[104]. The Emperor's New Mathematics: Western Learning and Imperial Authority%
\space%
During the Kangxi Reign (1662–1722). London: Oxford University Press, 2011.%
\par
\noindent\hangindent=2em\hangafter=1
% 原书页码：366；段落
[105] Jensen, Lionel. Manufacturing Confucianism: Chinese Traditions and Universal%
\space%
Civilization. Duke: Duke University Press, 1997.%
\par
\noindent\hangindent=2em\hangafter=1
% 原书页码：366；段落
[106]纪建勋，《〈碑记赠言合刻〉版本篇目互校》，自《国际汉学》，2014，第2期。%
\par
\noindent\hangindent=2em\hangafter=1
% 原书页码：366；段落
[107]蒋兆成，王日根，《康熙传》，北京，人民出版社，1998。%
\par
\noindent\hangindent=2em\hangafter=1
% 原书页码：366；段落
[108] Jin, Guoping. "The Mission to Congratulate the Kangxi Emperor on his Sev-%
entieth Birthday: Father João Mourão and His Plan to Bring an Ambassador of%
\space%
the King of Portugal to Beijing," in Macau Ricci Institute, eds., Acta Pekinensia:%
\space%
Western Historical Sources for the Kangxi Reign; International Symposium Orga-%
nized by the Macau Ricci Institute, Macau, 5th-7th October 2010. Macau: Macau%
\space%
Ricci Institute, 2013.%
\par
\noindent\hangindent=2em\hangafter=1
% 原书页码：366；段落
[109] . "Amicissimos' Tomas Pereira and Zhao Chang," in In the Light and Shad-%
ow of an Emperor. Tomás Pereira, S. J. (1645–1708), the Kangxi Emperor and the%
\space%
Jesuit Mission in China. Cambridge: Cambridge Scholar Publishing, 2012.%
\par
\noindent\hangindent=2em\hangafter=1
% 原书页码：366；段落
[110] Kato, Naoto. "Warrior Lamas: The Role of Lamas on Lobjang Danjin's Up-rising%
\space%
in Kokonor, 1723–1724," in Memoirs of the Research Department of Toyo Bunko%
\space%
62 (2004).%
\par
\noindent\hangindent=2em\hangafter=1
% 原书页码：366；段落
[111]. "Lobjang Danjin's Rebellion of 1723," in Gray Tuttle and Kurtis R. Schaef-%
fer, eds., The Tibetan History Reader. New York: Columbia University Press, 2013.%
\par
\noindent\hangindent=2em\hangafter=1
% 原书页码：367；段落
[112]鞠德源，《清宫廷画家郎世宁年谱》，自《故宫博物院院刊》，1988，第2期。%
\par
\noindent\hangindent=2em\hangafter=1
% 原书页码：367；段落
[113] Katz, Victor J., and Karen Hunger Parshall. Taming the Unknown: A History%
\space%
of Algebra from Antiquity to the Early Twentieth Century. Princeton: Princeton%
\space%
University Press, 2014.%
\par
\noindent\hangindent=2em\hangafter=1
% 原书页码：367；段落
[114] Kessler, Lawrence. K'ang-hsi and the Consolidation of Ch'ing Rule, 1661–1684.%
\space%
Chicago: University of Chicago Press, 1976.%
\par
\noindent\hangindent=2em\hangafter=1
% 原书页码：367；段落
[115] Kutcher, Norman A. Eunuch and Emperor in the Great Age of Qing Rule. Berke-%
ley: University of California Press, 2018.%
\par
\noindent\hangindent=2em\hangafter=1
% 原书页码：367；段落
[116] Lach, Donald. The Preface to Leibniz' Novissima Sinica: Commentary, Translation,%
\space%
Text. Honolulu: University of Hawaii Press, 1957.%
\par
\noindent\hangindent=2em\hangafter=1
% 原书页码：367；段落
[117] Landry-Deron, Isabelle. "The Kangxi Emperor's Lessons in Western Sciences as%
\space%
Recounted by the Jesuit Fathers J. Bouvet and J.-F. Gerbillon," in Macau Ricci%
\space%
Institute, eds., Acta Pekinensia: Western Historical Sources for the Kangxi Reign;%
\space%
International Symposium Organized by the Macau Ricci Institute, Macau, 5th-7th%
\space%
October 2010. Macau: Macau Ricci Institute, 2013.%
\par
\noindent\hangindent=2em\hangafter=1
% 原书页码：367；段落
[118] Larcombe, P. J. "The 18th century Chinese Discovery of the Catalan Numbers."%
Mathematical Spectrum 32, no.1 (1999).%
\par
\noindent\hangindent=2em\hangafter=1
% 原书页码：367；段落
[119] Latourette, Kenneth Scott. A History of Christian Missions in China. New York:%
\space%
The Macmillan Company, 1929.%
\par
\noindent\hangindent=2em\hangafter=1
% 原书页码：367；段落
[120] Lei, Zhang. "One City with Two Waters: Drinking Water in Beijing, 1644–%
1900," in Water History 9 (2017).%
\par
\noindent\hangindent=2em\hangafter=1
% 原书页码：367；段落
[121] Leslie, Donald. The Survival of the Chinese Jews: The Jewish Community of Kaifeng.%
\space%
Leiden, Brill, 1972.%
\par
\noindent\hangindent=2em\hangafter=1
% 原书页码：367；段落
[122] Li, Gertraude Roth. "The Rise of the Early Manchu State: A Portrait Drawn%
\space%
from Manchu Source to 1636." Ph.D. diss., Harvard University, 1975.%
\par
\noindent\hangindent=2em\hangafter=1
% 原书页码：367；段落
[123]. "State Building Before 1644," in The Cambridge History of China, Vol. 9.%
\space%
Cambridge: Cambridge University Press, 2002.%
\par
\noindent\hangindent=2em\hangafter=1
% 原书页码：367–368；段落（跨页续段）
[124]. "The Manchu-Chinese Relationship," in Jonathan Spence and John Wills,%
\space%
eds., From Ming to Ch'ing: Conquest, Region, and Continuity in Seventeenth-%
Century China. New Haven: Yale University Press, 1979.%
\par
\noindent\hangindent=2em\hangafter=1
% 原书页码：368；段落
[125]李天纲，《中国礼仪之争：历史、文献和意义》，上海：上海古籍出版社，%
1998。%
\par
\noindent\hangindent=2em\hangafter=1
% 原书页码：368；段落
[126] Lindorff, Joyce. "Missionaries, Keyboards and Musical Exchange in the Ming%
\space%
and Qing Courts," in Early Music 32, no. 3 (2004).%
\par
\noindent\hangindent=2em\hangafter=1
% 原书页码：368；段落
[127] Lippiello, Tiziana. "Astronomy and Astrology: Johann Adam Schall von Bell,"%
in Roman Malek, ed., Western Learning in China: The Contribution and Impact%
\space%
of Johann Adam Schall von Bell, S. J. (1592–1666). Nettetal: Steyler Verlag,%
\space%
1998.%
\par
\noindent\hangindent=2em\hangafter=1
% 原书页码：368；段落
[128]刘小萌，《八旗户籍中的旗下人诸名称考释》，自《社会科学辑刊》，%
1987，第3期。%
\par
\noindent\hangindent=2em\hangafter=1
% 原书页码：368；段落
[129] 鹿智钧，《根本与世仆：清朝旗人的法律地位》，台北：秀威出版，2017。%
\par
\noindent\hangindent=2em\hangafter=1
% 原书页码：368；段落
[130] Lui, Adam. Two Rulers in One Reign: Dorgon and Shun-Chih, 1644-1660.%
\space%
Australian National University Press, 1989.%
\par
\noindent\hangindent=2em\hangafter=1
% 原书页码：368；段落
[131]罗见今，《明安图是卡塔兰数的首创者》，《内蒙古大学学报（自然科学%
版）》，1988，第19卷，第2期。%
\par
\noindent\hangindent=2em\hangafter=1
% 原书页码：368；段落
[132] Malatesta, Edward. "A Fatal Clash of Wills: The Condemnation of the Chinese%
\space%
Rites by the Papal Legate Carlo Tommaso Maillard de Tournon," in David%
\space%
Mungello, ed., The Chinese Rites Controversy: Its History and Meaning. San%
\space%
Francisco: The Ricci Institute for Chinese-Western Cultural History, 1994.%
\par
\noindent\hangindent=2em\hangafter=1
% 原书页码：368；段落
[133] Malek, Roman, ed., Western Learning and Christianity in China: The Contri-%
bution and Impact of Johann Adam Schall von Bell, S. J. (1592–1666). Rout-%
ledge, 1998.%
\par
\noindent\hangindent=2em\hangafter=1
% 原书页码：368；段落
[134] Marinescu, Jocelyn M. N. "Defending Christianity in China: The Lettres%
édifiantes et curieuses et Ruijianlu in Relation to the Yongzheng Proscription of%
\space%
1724." Ph.D. diss., Kansas State University, 2008.%
\par
\noindent\hangindent=2em\hangafter=1
% 原书页码：368–369；段落（跨页续段）
[135] Menegon, Eugenio. Ancestors, Virgins, and Friars: Christianity as a Local Religion in%
\space%
Late Imperial China. Cambridge: Harvard University Press, 2009.%
\par
\noindent\hangindent=2em\hangafter=1
% 原书页码：369；段落
[136]. "Yongzheng's Conundrum. The Emperor on Christianity, Religions, and%
\space%
Heterodoxy," in Hoster, Barbara, Kuhlmann, Dirk and Wesolowski, Zbigniew,%
\space%
eds., Rooted in Hope / In Der Hoffnung Verwurzelt. Festschrift in Honor of Roman%
\space%
Malek S. V. D. on the Occasion of His 65th Birthday. Routledge: Monumenta Seri-%
ca Institute I (2017).%
\par
\noindent\hangindent=2em\hangafter=1
% 原书页码：369；段落
[137]. "Interlopers at the Fringes of Empire: The Procurators of the Propaganda%
\space%
Fide Papal Congregation in Canton and Macau, 1700-–1823," in Cross-Currents:%
\space%
East Asian History and Culture Review 7 (2018).%
\par
\noindent\hangindent=2em\hangafter=1
% 原书页码：369；段落
[138]. "A Clash of Court Cultures: Papal Envoys in Early Eighteenth Century%
\space%
Beijing," in Luís Filipe Barreto, ed., Europe-China. Intercultural Encounters%
(16th-18th Centuries). Lisbon: Centro Científico e Cultural de Macau,%
\space%
2012.%
\par
\noindent\hangindent=2em\hangafter=1
% 原书页码：369；段落
[139]孟森，《八旗制度考实》，自《明清史论著集刊》，北京：中华书局，%
1959。%
\par
\noindent\hangindent=2em\hangafter=1
% 原书页码：369；段落
[140]孟森，《清世宗人承大统考实》，自《明清史论著集刊》，北京：中华书局，%
1959。%
\par
\noindent\hangindent=2em\hangafter=1
% 原书页码：369；段落
[141]孟森，《清代史》，台北：正中书局，1962。%
\par
\noindent\hangindent=2em\hangafter=1
% 原书页码：369；段落
[142]孟森，《清世祖董鄂妃生死特殊典礼》，自《明清史论著集刊续编》，北京：%
中华书局，1986。%
\par
\noindent\hangindent=2em\hangafter=1
% 原书页码：369；段落
[143]孟森，《世祖出家事考实》，自《明清史论著集刊续编》，北京：中华书局，%
1986。%
\par
\noindent\hangindent=2em\hangafter=1
% 原书页码：369；段落
[144]孟森，《明元清系通纪》，台北：广文书局，1972，1934年初版。%
\par
\noindent\hangindent=2em\hangafter=1
% 原书页码：369；段落
[145]孟森，《清初三大疑案考实》，成都：巴蜀书社，2002。%
\par
\noindent\hangindent=2em\hangafter=1
% 原书页码：369；段落
[146] 孟昭信，《康熙大帝全传》，长春：吉林文史出版社，1987。%
\par
\noindent\hangindent=2em\hangafter=1
% 原书页码：369–370；段落（跨页续段）
[147] Miasnikov, Vladimir. "Ferdinand Verbiest and His Role in the Formation of%
\space%
Sino-Russian Diplomatic Relations," in John Witek, ed., Ferdinand Verbiest%
(1623–1688): Jesuit Missionary, Scientist, Engineer and Diplomat. Nettetal: Steyler%
\space%
Verlag, 1994.%
\par
\noindent\hangindent=2em\hangafter=1
% 原书页码：370；段落
[148] Minamiki, George. The Chinese Rites Controversy from Its Beginning to Modern%
\space%
Times. Chicago: Loyola University Press, 1985.%
\par
\noindent\hangindent=2em\hangafter=1
% 原书页码：370；段落
[149]牟润孙，《注史斋丛稿》，北京：中华书局，1987。%
\par
\noindent\hangindent=2em\hangafter=1
% 原书页码：370；段落
[150] Mungello, David. Curious Land: Jesuit Accommodation and the Origins of Sinology.%
\space%
Wiesbaden: Steiner Verlag, 1985.%
\par
\noindent\hangindent=2em\hangafter=1
% 原书页码：370；段落
[151] Musillo, Marco. "Reconciling Two Careers: The Jesuit Memoir of Giuseppe Cas-%
tiglione Lay Brother and Qing Imperial Painter," in Eighteenth-Century Studies%
\space%
42, no. 1 (2008).%
\par
\noindent\hangindent=2em\hangafter=1
% 原书页码：370；段落
[152]. Bridging Europe and China: The Professional life of Giuseppe Castiglione%
(1688–1766). Ph. D. diss., University of East Anglia, 2006.%
\par
\noindent\hangindent=2em\hangafter=1
% 原书页码：370；段落
[153] Naquin, Susan. "Giuseppe Castiglione/Lang Shining: A Review Essay," in T'oung%
\space%
Pao, Second Series, 95, no. 4/5 (2009).%
\par
\noindent\hangindent=2em\hangafter=1
% 原书页码：370；段落
[154] Oxnam, Robert B, Ruling from Horseback: Manchu Politics in the Oboi Regency,%
\space%
1661–1669. Chicago and London: University of Chicago Press, 1975.%
\par
\noindent\hangindent=2em\hangafter=1
% 原书页码：370；段落
[155]. "Policies and Institutions of the Oboi Regency, 1661–1669," in The Jour-%
nal of Asian Studies 32, no. 2 (1973).%
\par
\noindent\hangindent=2em\hangafter=1
% 原书页码：370；段落
[156] Pagani, Catherine. "Eastern Magnificence \& European Ingenuity": Clocks of Late%
\space%
Imperial China. Ann Arbor: University of Michigan Press, 2001.%
\par
\noindent\hangindent=2em\hangafter=1
% 原书页码：370；段落
[157] Pelliot, Paul. "Michel Boym," in T’oung Pao, Second Series, 31, no. 1/2 (1934).%
\par
\noindent\hangindent=2em\hangafter=1
% 原书页码：370；段落
[158] Perdue, Peter. China Marches West: The Qing Conquest of Central Eurasia. Cam-%
bridge: Harvard University Press, 2005.%
\par
\noindent\hangindent=2em\hangafter=1
% 原书页码：370；段落
[159] Petech, Luciano. China and Tibet in the Early 18th Century: History of the Estab-%
lishment of Chinese Protectorate in Tibet. Leiden: Brill, 1950.%
\par
\noindent\hangindent=2em\hangafter=1
% 原书页码：370；段落
[160] Piastra, Stefano. "The Shenjianlu by Fan Shouyi Reconsidered: New Geo-Histori-%
cal Notes," in Fudan Journal of the Humanities and Social Sciences 5, no. 4 (2012).%
\par
\noindent\hangindent=2em\hangafter=1
% 原书页码：371；段落
[161] Pih, Irene. Le père Gabriel de Magalhães. Un jésuite portugais en Chine au XVIIe%
\space%
siècle. Paris: Fundação Calouste Gulbenkian-Centro Cultural Português, 1979.%
\par
\noindent\hangindent=2em\hangafter=1
% 原书页码：371；段落
[162]祁美琴，崔灿，《包衣身份再辨》，自《清史研究》，2013，第1期。%
\par
\noindent\hangindent=2em\hangafter=1
% 原书页码：371；段落
[163]祁美琴，《清初内务府及其与十三衙门的关系》，自《清史研究》，1997，%
第1期。%
\par
\noindent\hangindent=2em\hangafter=1
% 原书页码：371；段落
[164]祁美琴，《清代内务府》，北京：中国人民大学出版社，1998。%
\par
\noindent\hangindent=2em\hangafter=1
% 原书页码：371；段落
[165] Rawski, Evelyn. "Ch'ing Imperial Marriage and Problems of Rulership," in Rubie%
\space%
Waston and Patricia Buckly Ebrey, eds., Marriage and Inequality in Chinese%
\space%
Society. Berkeley: University California Press, 1991.%
\par
\noindent\hangindent=2em\hangafter=1
% 原书页码：371；段落
[166]. "Presidential Address: Reenvisioning the Qing: The Significance of the%
\space%
Qing Period in Chinese History," in Journal of Asian Studies 55, no. 4 (1996).%
\par
\noindent\hangindent=2em\hangafter=1
% 原书页码：371；段落
[167]. The Last Emperors: A Social History of Qing Imperial Institutions. Berkeley:%
\space%
University of California Press, 1998.%
\par
\noindent\hangindent=2em\hangafter=1
% 原书页码：371；段落
[168] Rinaldi, Bianca Maria. Ideas of Chinese Gardens: Western Accounts, 1300–1860.%
\space%
Philadelphia: University of Pennsylvania Press, 2016.%
\par
\noindent\hangindent=2em\hangafter=1
% 原书页码：371；段落
[169] Rockhill, W. W. "The Dalai Lamas of Lhasa and Their Relations with the Man-%
chu Emperors of China, 1644–1908," in T’oung Pao, Second Series, 11, no. 1%
(1910).%
\par
\noindent\hangindent=2em\hangafter=1
% 原书页码：371；段落
[170] Rosso, Antonio Sisto. Apostolic Legations to China of the Eighteenth Century.%
\space%
South Pasadena: P. D. and I. Perkins, 1948.%
\par
\noindent\hangindent=2em\hangafter=1
% 原书页码：371；段落
[171] Rouleau, Francis. "Maillard de Tournon, Papal legate at the court of Peking: the%
\space%
first imperial audience (31 December 1705)," in Archivum historicum Societatis%
\space%
Iesu 31 (1962).%
\par
\noindent\hangindent=2em\hangafter=1
% 原书页码：371；段落
[172] . "The First Chinese Priest of the Society of Jesus: Emmanuel de Siqueira.%
\space%
Cheng-Ma-no Wei-hsin. 1633-1673," in Archivum Historicum Societatis Iesu 28%
(1959).%
\par
\noindent\hangindent=2em\hangafter=1
% 原书页码：371–372；段落（跨页续段）
[173] Rowbotham, Arnold. Missionary and Mandarin: The Jesuits at the Court of China.%
\space%
Berkeley: University of California Press, 1942.%
\par
\noindent\hangindent=2em\hangafter=1
% 原书页码：372；段落
[174] Rule, Paul. "Louis Fan Shouyi and Macau," in Review of Culture 21 (1994).%
\par
\noindent\hangindent=2em\hangafter=1
% 原书页码：372；段落
[175]. Kung-Tsu or Confucius: The Jesuit Interpretation of Confucianism. Boston:%
\space%
Allen \& Unwin Press, 1986.%
\par
\noindent\hangindent=2em\hangafter=1
% 原书页码：372；段落
[176] . "Louis Fan Shou-i: a Missing Link in the Chinese Rites Controversy," in%
\space%
Edward Malatesta, ed., Échanges culturels et religieux entre la Chine et I'Occident.%
\space%
San Francisco: The Ricci Institute for Chinese-Western Cultural History, 1995.%
\par
\noindent\hangindent=2em\hangafter=1
% 原书页码：372；段落
[177] Schwieger, Peter. The Dalai Lama and the Emperor of China: A Political History%
\space%
of the Tibetan Institution of Reincarnation. New York: Columbia University Press,%
\space%
2014.%
\par
\noindent\hangindent=2em\hangafter=1
% 原书页码：372；段落
[178] Sebes, Joseph. The Jesuits and the Sino-Russian Treaty of Nerchinsk (1689): The Di-%
ary of Thomas Pereira, S. J.. Rome: Institutum Historicum, 1961.%
\par
\noindent\hangindent=2em\hangafter=1
% 原书页码：372；段落
[179] Sen, Sudipta. "The New Frontier of Manchu China and the Historiography of%
\space%
Asian Empire: A Review Essay," in Journal of Asian Studies 61, no. 1 (2002).%
\par
\noindent\hangindent=2em\hangafter=1
% 原书页码：372；段落
[180] Sewell, John. Elementary Treatise on Steam and Locomotion. London: John Weale,%
\space%
1852.%
\par
\noindent\hangindent=2em\hangafter=1
% 原书页码：372；段落
[181]史松，《康熙朝皇位继承斗争和雍正继位》，自中国人民大学清史研究所%
编，《清史研究集（第四辑）》，成都：四川人民出版社，1986。%
\par
\noindent\hangindent=2em\hangafter=1
% 原书页码：372；段落
[182]圣空法师，《清世宗与佛教》，中华佛学研究所，硕士论文，2000。%
\par
\noindent\hangindent=2em\hangafter=1
% 原书页码：372；段落
[183] Shu, Liguang. "Ferdinand Verbiest and the Casting of Cannons in the Qing%
\space%
Dynasty," in John Witek, ed., Ferdinand Verbiest (1623–1688): Jesuit Missionary,%
\space%
Scientist, Engineer and Diplomat. Nettetal: Steyler Verlag, 1994.%
\par
\noindent\hangindent=2em\hangafter=1
% 原书页码：372；段落
[184]舒理广等，《南怀仁与中国清代铸造的大炮》，自《故宫博物院院刊》，%
1989，第1期。%
\par
\noindent\hangindent=2em\hangafter=1
% 原书页码：372；段落
[185] Smith, Richard J. "Divination in Ching Dynasty China," in Richard J. Smith%
\space%
and D. W. Y. Kwok, eds., Cosmology, Ontology, and Human Efficacy: Essays in%
\space%
Chinese Thought. Honolulu: University of Hawaii Press, 1993.%
\par
\noindent\hangindent=2em\hangafter=1
% 原书页码：373；段落
[186]. Fortune-tellers and Philosophers: Divinations in Traditional Chinese Society.%
\space%
Boulder: Westview Press, 1991.%
\par
\noindent\hangindent=2em\hangafter=1
% 原书页码：373；段落
[187] Soloshcheva, Maria. "The Conquest of Qinghai' Stele of 1725 and the After-%
math of Lobsang Danjin's Rebellion in 1723–1724," in Higher School of Econom-%
ics Research Paper, 2014.%
\par
\noindent\hangindent=2em\hangafter=1
% 原书页码：373；段落
[188] Sommer, Matthew. Sex, Law, and Society in Late Imperial China. Stanford:%
\space%
Stanford University Press, 2000.%
\par
\noindent\hangindent=2em\hangafter=1
% 原书页码：373；段落
[189]宋秉仁，《从文臣谥号看明清谥法理论》，自《白沙人文社会学报》，%
2002，第1期。%
\par
\noindent\hangindent=2em\hangafter=1
% 原书页码：373；段落
[190] Spence, Jonathan. "Claims and Counter-claims: The Kangxi Emperor and the%
\space%
Europeans (1661–1722)," in David Mungello, ed., The Chinese Rites Controversy:%
\space%
Its History and Meaning. San Francisco: The Ricci Institute for Chinese-Western%
\space%
Cultural History, 1994.%
\par
\noindent\hangindent=2em\hangafter=1
% 原书页码：373；段落
[191] . Emperor of China: Self-Portrait of Kang-hsi. New York: Vintage Books,%
\space%
1974.%
\par
\noindent\hangindent=2em\hangafter=1
% 原书页码：373；段落
[192] . "The Kang-hsi Reign," in The Cambridge History of China: The Ch'ing%
\space%
Empire to 1800, Vol. 9. Cambridge: Cambridge University Press, 2002.%
\par
\noindent\hangindent=2em\hangafter=1
% 原书页码：373；段落
[193]. The Search for Modern China. New York: W. W. Norton, 1st edition, 1991.%
\par
\noindent\hangindent=2em\hangafter=1
% 原书页码：373；段落
[194]. Treason by the Book. New York: Viking, 2001.%
\par
\noindent\hangindent=2em\hangafter=1
% 原书页码：373；段落
[195]. T’ao Yin and the K'ang-hsi Emperor: Bondservant and Master. New Haven:%
\space%
Yale University Press, 1966.%
\par
\noindent\hangindent=2em\hangafter=1
% 原书页码：373；段落
[196] Standaert, Nicolas. "The Edict of Tolerance: A Textual History and Reading," in%
\space%
Wardega and de Saldanhapp, eds., In the Light and Shadow of an Emperor. Tomás%
\space%
Pereira, S. J. (1645–1708), the Kangxi Emperor and the Jesuit Mission in China.%
\space%
Cambridge: Cambridge Scholar Publishing, 2012.%
\par
\noindent\hangindent=2em\hangafter=1
% 原书页码：373；段落
[197]  ed., Handbook of Christianity in China: Volume One, 635-1800. Brill:%
\space%
Leiden, Boston, Koln, 2001.%
\par
\noindent\hangindent=2em\hangafter=1
% 原书页码：373–374；段落（跨页续段）
[198]汤开建，《雍正教难期间驱逐传教士至广州事件始末考》，《清史研究》，%
2014，第2期。%
\par
\noindent\hangindent=2em\hangafter=1
% 原书页码：374；段落
[199]汤开建，《沉与浮：明清鼎革变局中的欧洲传教士利类思与安文思》（上、%
下），《北京行政学院学报》，2014年第4、5期。%
\par
\noindent\hangindent=2em\hangafter=1
% 原书页码：374；段落
[200] Tarsetti, Gabriele, "Mariner, musician and missionary, and true priest always:%
\space%
Teodorico Pedrini's life in Xitang," in Proceedings of the papers on Teodorico Pedrini%
\space%
for the 11th Symposium on the History of the Church in China – KU Leuven 4th-%
7th September 2012.%
\par
\noindent\hangindent=2em\hangafter=1
% 原书页码：374；段落
[201] Torbert, Preston. The Ch’ing Imperial Household Department: A Study of Its%
\space%
Organization and Principal Functions, 1662-1796. Cambridge: Harvard%
\space%
University Press, 1977.%
\par
\noindent\hangindent=2em\hangafter=1
% 原书页码：374；段落
[202] Tuttle, Gray. "Tibetan Buddhism at Wutai Shan in the Qing," in The Journal of%
\space%
the International Association of Tibetan Studies 6 (2011).%
\par
\noindent\hangindent=2em\hangafter=1
% 原书页码：374；段落
[203] Wakeman, Frederic. The Great Enterprise: the Manchu Reconstruction of the%
\space%
Imperial Order in Seventeenth-Century China. Berkeley: University of California%
\space%
Press, 1985.%
\par
\noindent\hangindent=2em\hangafter=1
% 原书页码：374；段落
[204] Waley-Cohen, Joanna. "The New Qing History," in Radical History Review 88%
(2004).%
\par
\noindent\hangindent=2em\hangafter=1
% 原书页码：374；段落
[205] Walle, Willy Vande. "Ferdinand Verbiest and the Chinese Bureaucracy," in John%
\space%
Witek, ed., Ferdinand Verbiest (1623–1688): Jesuit Missionary, Scientist, Engineer%
\space%
and Diplomat. Nettetal: Steyler Verlag, 1994.%
\par
\noindent\hangindent=2em\hangafter=1
% 原书页码：374；段落
[206]王冕森，《允禟西洋密码字档案初探》，《清史研究》，2021，第5期。%
\par
\noindent\hangindent=2em\hangafter=1
% 原书页码：374；段落
[207] Wang, Xiangyun. "The Qing Court's Tibet Connection: Lcang skya Rol pa'i rdo%
\space%
rje and the Qianlong Emperor," in Harvard Journal of Asiatic Studies 60, no. 1%
(2000).%
\par
\noindent\hangindent=2em\hangafter=1
% 原书页码：374；段落
[208]王钟翰，《清世宗夺嫡考实》，自王钟翰，《清史杂考》，北京：中华书局，%
1963。%
\par
\noindent\hangindent=2em\hangafter=1
% 原书页码：374；段落
[209] Wills, John. Embassies and Illusions: Dutch and Portuguese Envoys to K'ang-hsi,%
\space%
1666–1687. Boston: Harvard University Press, 1984.%
\par
\noindent\hangindent=2em\hangafter=1
% 原书页码：375；段落
[210] Witek, John. Controversial Ideas in China and in Europe: a Biography of Jean-%
François Foucquet, S. J., (1665–1741). Roma: Institutum Historicum S. I., 1982.%
\par
\noindent\hangindent=2em\hangafter=1
% 原书页码：375；段落
[211] . "Johann Adam Schall von Bell and the Transition from the Ming to the%
\space%
Ch'ing Dynasty," in R. Malek, ed., Western Learning and Christianity in China.%
\space%
The Contribution and Impact of Johann Adam Schall von Bell, S. J. (1592–1666).%
\space%
Nettetal: Steyler Verlag, 1998.%
\par
\noindent\hangindent=2em\hangafter=1
% 原书页码：375；段落
[212]. "Understanding the Chinese: a Comparison of Matteo Ricci and the French%
\space%
Jesuit Mathematicians Sent by Louis XIV, " in Charles E. Ronan and Bonnie B. C.,%
\space%
eds., East Meets West: The Jesuits in China, 1582–1773. Chicago: Loyola University%
\space%
Press, 1988.%
\par
\noindent\hangindent=2em\hangafter=1
% 原书页码：375；段落
[213]. “Eliminating Misunderstandings: Antoine de Beauvollier (1657–1708) and%
\space%
His Eclaircissements Sur Les Controversies De La Chine," in David Mungello,%
\space%
ed., The Chinese Rites Controversy: Its History and Meaning. San Francisco: The%
\space%
Ricci Institute for Chinese-Western Cultural History, 1994.%
\par
\noindent\hangindent=2em\hangafter=1
% 原书页码：375；段落
[214]魏特著，杨丙辰译，《汤若望传》，台北：台湾商务印书馆，1960。%
\par
\noindent\hangindent=2em\hangafter=1
% 原书页码：375；段落
[215] 吴伯娅，《关于雍正禁教的几个问题：耶稣会士书简与清代档案的比读》%
自《清史论丛2003—2004年号》，北京：中国广播电视出版社，2004。%
\par
\noindent\hangindent=2em\hangafter=1
% 原书页码：375；段落
[216] Wu, Silas H. L. Communication and Imperial Control in China: Evolution of the%
\space%
Palace Memorial System. Cambridge: Harvard University Press, 1970.%
\par
\noindent\hangindent=2em\hangafter=1
% 原书页码：375；段落
[217]. Passage to Power: K'ang-hsi and His Heir Apparent, 1661–1722. Cambridge:%
\space%
Harvard University Press, 1979.%
\par\noindent\hangindent=2em\hangafter=1
[218] 夏伯嘉，《从安文思与汤若望之间的恩怨论说明宫廷耶稣会士》，《北京%
行政学院学报》，2017，第3期.%
\par
\noindent\hangindent=2em\hangafter=1
% 原书页码：375；段落
[219] 萧若瑟，《天主教传行中国考》，上海：上海书店，1989。%
\par
\noindent\hangindent=2em\hangafter=1
% 原书页码：375；段落
[220]谢景芳，《天花与清初史实评议》，自《民族研究》，1995，第1期。%
\par
\noindent\hangindent=2em\hangafter=1
% 原书页码：375；段落
[221]虚云，《上海玉佛寺禅七开示》，自《虚云法师讲演录》，北京：中华书局，%
2018。%
\par
\noindent\hangindent=2em\hangafter=1
% 原书页码：375；段落
[222]徐宗泽，《明清间耶稣会士译著提要》，台北：台湾中华书局，1958。%
\par
\noindent\hangindent=2em\hangafter=1
% 原书页码：375；段落
[223]徐宗泽，《中国天主教传教史概论》，上海：上海书店出版社，1990。%
\par
\noindent\hangindent=2em\hangafter=1
% 原书页码：376；段落
[224]杨珍，《康熙朝鳌拜罪案辨析》，自《历史档案》，2016，第3期。%
\par
\noindent\hangindent=2em\hangafter=1
% 原书页码：376；段落
[225]印光，《拣魔辨异录石印序》，自《印光法师文钞（第三卷）》，北京：宗%
教文化出版社，2000。%
\par
\noindent\hangindent=2em\hangafter=1
% 原书页码：376；段落
[226] Young, John D. "An Early Confucian Attack on Christianity: Yang Kuang-hsien%
\space%
and His Pu-te-I," in Journal of Chinese University of Hongkong 3 (1975).%
\par
\noindent\hangindent=2em\hangafter=1
% 原书页码：376；段落
[227] . Confucianism and Christianity: The First Encounter. Hong Kong: Hong%
\space%
Kong University Press, 1983.%
\par
\noindent\hangindent=2em\hangafter=1
% 原书页码：376；段落
[228] Zelin, Madeleine. "Yung-Cheng Reign," in The Cambridge History of China: The%
\space%
Ch’ing Empire to 1800. Vol. 9. Cambridge: Cambridge University Press, 2002.%
\par
\noindent\hangindent=2em\hangafter=1
% 原书页码：376；段落
[229]张小青，《明清之际西洋火炮的输人及其影响》，自《清史研究集（第四%
辑）》，成都：四川人民出版社，1986。%
\par
\noindent\hangindent=2em\hangafter=1
% 原书页码：376；段落
[230]郑天挺，《清史探微》，北京：北京大学出版社，1999。%
\par
\noindent\hangindent=2em\hangafter=1
% 原书页码：376；段落
[231]郑天挺，《清代包衣制度与宦官》，自《探微集》，北京：中华书局，1980。%
\par
\noindent\hangindent=2em\hangafter=1
% 原书页码：376；段落
[232] 周叔迦，《清代佛教史料辑稿》，台北：新文丰出版公司，2000。%
\par
\noindent\hangindent=2em\hangafter=1
% 原书页码：376；段落
[233]朱维铮，《走出中世纪》，上海：上海人民出版社，1987。%
\par
\noindent\hangindent=2em\hangafter=1
% 原书页码：376；段落
[234]左云鹏，《清代旗下奴仆的地位及其变化》，自《陕西师范大学学报（哲%
学社会科学版）》，1980，第1期。%
\par
\noindent\hangindent=2em\hangafter=1
% 原书页码：376；段落
[235] Zürcher, Erik. "In the Yellow Tiger's Den: Buglio and Magalháes at the Court of%
\space%
Zhang Xianzhong, 1644–1647," in Monumenta Serica 50 (2002).%
\par
\noindent\hangindent=2em\hangafter=1
% 原书页码：376；段落
[236] . “Jesuit Accommodation and the Chinese Cultural Imperative," in David%
\space%
Mungello, ed., The Chinese Rites Controversy: Its History and Meaning. San%
\space%
Francisco: The Ricci Institute for Chinese-Western Cultural History, 1994.%
\par
\noindent\hangindent=2em\hangafter=1
% 原书页码：376；段落
[237] Zurndorfer, Harriet T. "One Adam Having Driven Us out of Paradise, Another%
\space%
Has Driven Us out of China': Yang Kuang-hsien's Challenge of Adam Schall%
\space%
von Bell," in Leonard Blussé and Harriet T. Zurndorfer, eds., Conflict and%
\space%
Accommodation in Early Modern China: Essays in Honour of Erik Zürcher. Leiden:%
\space%
Brill, 1993.%
